% Options for packages loaded elsewhere
\PassOptionsToPackage{unicode}{hyperref}
\PassOptionsToPackage{hyphens}{url}
\documentclass[
  russian,
  11pt,
]{article}
\usepackage{xcolor}
\usepackage[margin=2.5cm]{geometry}
\usepackage{amsmath,amssymb}
\setcounter{secnumdepth}{-\maxdimen} % remove section numbering
\usepackage{iftex}
\ifPDFTeX
  \usepackage[T1]{fontenc}
  \usepackage[utf8]{inputenc}
  \usepackage{textcomp} % provide euro and other symbols
\else % if luatex or xetex
  \usepackage{unicode-math} % this also loads fontspec
  \defaultfontfeatures{Scale=MatchLowercase}
  \defaultfontfeatures[\rmfamily]{Ligatures=TeX,Scale=1}
\fi
\usepackage{lmodern}
\ifPDFTeX\else
  % xetex/luatex font selection
\fi
% Use upquote if available, for straight quotes in verbatim environments
\IfFileExists{upquote.sty}{\usepackage{upquote}}{}
\IfFileExists{microtype.sty}{% use microtype if available
  \usepackage[]{microtype}
  \UseMicrotypeSet[protrusion]{basicmath} % disable protrusion for tt fonts
}{}
\makeatletter
\@ifundefined{KOMAClassName}{% if non-KOMA class
  \IfFileExists{parskip.sty}{%
    \usepackage{parskip}
  }{% else
    \setlength{\parindent}{0pt}
    \setlength{\parskip}{6pt plus 2pt minus 1pt}}
}{% if KOMA class
  \KOMAoptions{parskip=half}}
\makeatother
% definitions for citeproc citations
\NewDocumentCommand\citeproctext{}{}
\NewDocumentCommand\citeproc{mm}{%
  \begingroup\def\citeproctext{#2}\cite{#1}\endgroup}
\makeatletter
 % allow citations to break across lines
 \let\@cite@ofmt\@firstofone
 % avoid brackets around text for \cite:
 \def\@biblabel#1{}
 \def\@cite#1#2{{#1\if@tempswa , #2\fi}}
\makeatother
\newlength{\cslhangindent}
\setlength{\cslhangindent}{1.5em}
\newlength{\csllabelwidth}
\setlength{\csllabelwidth}{3em}
\newenvironment{CSLReferences}[2] % #1 hanging-indent, #2 entry-spacing
 {\begin{list}{}{%
  \setlength{\itemindent}{0pt}
  \setlength{\leftmargin}{0pt}
  \setlength{\parsep}{0pt}
  % turn on hanging indent if param 1 is 1
  \ifodd #1
   \setlength{\leftmargin}{\cslhangindent}
   \setlength{\itemindent}{-1\cslhangindent}
  \fi
  % set entry spacing
  \setlength{\itemsep}{#2\baselineskip}}}
 {\end{list}}
\usepackage{calc}
\newcommand{\CSLBlock}[1]{\hfill\break\parbox[t]{\linewidth}{\strut\ignorespaces#1\strut}}
\newcommand{\CSLLeftMargin}[1]{\parbox[t]{\csllabelwidth}{\strut#1\strut}}
\newcommand{\CSLRightInline}[1]{\parbox[t]{\linewidth - \csllabelwidth}{\strut#1\strut}}
\newcommand{\CSLIndent}[1]{\hspace{\cslhangindent}#1}
\ifLuaTeX
\usepackage[bidi=basic,shorthands=off,provide=*]{babel}
\else
\usepackage[bidi=default,shorthands=off,provide=*]{babel}
\fi
\ifLuaTeX
  \usepackage{selnolig} % disable illegal ligatures
\fi
\setlength{\emergencystretch}{3em} % prevent overfull lines
\providecommand{\tightlist}{%
  \setlength{\itemsep}{0pt}\setlength{\parskip}{0pt}}
% Преамбула для сборки статьи (xelatex): кириллица, booktabs, математика.
% babel здесь НЕ подключается: его добавляет шаблон pandoc по метаданным lang.
\usepackage{fontspec}
\setmainfont{Liberation Serif}
\setsansfont{Liberation Sans}
\setmonofont{Liberation Mono}[Scale=0.86]

\usepackage{booktabs}
\usepackage{longtable}
\usepackage{array}
\usepackage{amsmath,amssymb}
\usepackage{caption}
\usepackage{float}
\usepackage{graphicx}

\usepackage[unicode,colorlinks=true,linkcolor=blue!45!black,citecolor=blue!45!black,urlcolor=blue!45!black]{hyperref}

\captionsetup{font=small,labelfont=bf,justification=justified,singlelinecheck=false}
\setlength{\parskip}{0.35em}
\setlength{\parindent}{0pt}

% Русские подписи («Таблица», «Рисунок», «Литература») даёт babel-russian
% автоматически, поскольку в метаданных указано lang: ru-RU.
\usepackage{bookmark}
\IfFileExists{xurl.sty}{\usepackage{xurl}}{} % add URL line breaks if available
\urlstyle{same}
\hypersetup{
  pdftitle={Льготная ипотека как автономный источник спроса и трансмиссия денежно-кредитной политики: свидетельства разностных разностей по закрытию массовой программы},
  pdfauthor={Ленар (исследовательский трек, профиль research)},
  pdflang={ru-RU},
  pdfkeywords={льготная ипотека; трансмиссия денежно-кредитной политики;
рынок жилья; разностные разности; рандомизационный вывод; региональная
гетерогенность},
  hidelinks,
  pdfcreator={LaTeX via pandoc}}

\title{Льготная ипотека как автономный источник спроса и трансмиссия
денежно-кредитной политики: свидетельства разностных разностей по
закрытию массовой программы}
\author{Ленар (исследовательский трек, профиль research)}
\date{1 октября 2026 г.}

\begin{document}
\maketitle
\begin{abstract}
Работа отвечает на вопрос, действительно ли программы льготной ипотеки
формировали автономный спрос на жильё, ослабляя передачу ключевой
ставки, и что происходит с выдачами, когда такая программа закрывается.
Идентификация построена на закрытии массовой программы «Льготная ипотека
8\%» 1 июля 2024 года и непрерывной мере экспозиции --- доле выдач по
программе в общих жилищных выдачах региона за двенадцать месяцев до
события. На панели из 84 субъектов Российской Федерации (июль 2023 ---
июнь 2025) с фиксированными эффектами региона и месяца, региональными
линейными трендами и кластеризованными по региону стандартными ошибками
оценка эффекта на объём выдач составляет −0,557 (t = −3,84);
рандомизационный тест по 2 000 перестановок экспозиции даёт p = 0,003.
Событийное исследование показывает незначимые предсобытийные
коэффициенты (\textbar t\textbar{} ≤ 1,61) и растущий по модулю
постсобытийный эффект: от −0,518 на квартале события до −1,593 на
третьем квартале после него, что при медианной экспозиции соответствует
падению объёма на 22,6 \%, а в Ленинградской области --- на 54,2 \%.
Плацебо-событие 1 июля 2023 года эффект не воспроизводит. Эффект выражен
в объёме, но не в числе кредитов, то есть программа поддерживала прежде
всего размер кредита, а не число сделок. На ограниченной уральской
панели (девять субъектов) тот же дизайн не идентифицируется: знак
коэффициента меняется от спецификации к спецификации, рандомизационный
тест даёт p = 0,655, а плацебо-событие значимо. Этот контраст
показывает, что для оценки дифференциальных эффектов жилищной политики
решающее значение имеет широта контрольной группы, а не набор контролей.
\end{abstract}

\textbf{Ключевые слова:} льготная ипотека; трансмиссия денежно-кредитной
политики; рынок жилья; разностные разности; рандомизационный вывод;
региональная гетерогенность.

\textbf{Классификация JEL:} E52, G21, R31, R38.

\section{Abstract (English)}\label{abstract-english}

This paper asks whether subsidised mortgage programmes in Russia
generated an autonomous source of housing demand, weakening the
transmission of the policy rate, and what happens to lending when such a
programme is terminated. Identification exploits the closure of the mass
programme ``Mortgage at 8\%'\,' on 1 July 2024 and a continuous measure
of regional exposure --- the share of programme loans in total housing
loans over the twelve months preceding the event. On a panel of 84
Russian regions (July 2023 to June 2025) with region and month fixed
effects, region-specific linear trends and region-clustered standard
errors, the effect on lending volume is \(-0{,}557\) (\(t = -3{,}84\));
a randomisation test over 2\{,\}000 reassignments of exposure yields
\(p = 0{,}003\). An event study reports insignificant pre-event
coefficients (\(|t| \le 1{,}61\)) and a monotonically growing post-event
effect, from \(-0{,}518\) in the quarter of the event to \(-1{,}593\)
three quarters later --- equivalent to a 22.6\% decline in volume at
median exposure and 54.2\% in Leningrad Oblast. A placebo event on 1
July 2023 does not reproduce the result. The effect operates through
loan volume rather than loan counts, indicating that the subsidy
supported the size of loans rather than the number of transactions. On a
restricted panel of nine Urals regions the same design fails to identify
the effect: the coefficient changes sign across specifications, the
randomisation test gives \(p = 0{,}655\), and the placebo event is
significant. The contrast shows that the breadth of the control group,
not the set of covariates, is decisive for identification in this class
of design.

\textbf{Keywords:} subsidised mortgages; monetary policy transmission;
housing market; difference-in-differences; randomisation inference;
regional heterogeneity.

\textbf{JEL classification:} E52, G21, R31, R38.

\section{Введение}\label{ux432ux432ux435ux434ux435ux43dux438ux435}

Российский рынок жилья последние пять лет работает в режиме, который
плохо описывается стандартной моделью трансмиссии денежно-кредитной
политики (ДКП). Ключевая ставка выросла с 7,5 \% в середине 2023 года до
21 \% к концу 2024 года, однако выдачи ипотечных кредитов в значительной
части регионов не следовали за этим движением с ожидаемой силой:
значительная доля сделок шла по программам с субсидируемой ставкой,
слабо связанной с ключевой. Возникает вопрос, который и составляет
предмет настоящей работы: формировали ли программы льготной ипотеки
автономный, то есть независимый от ставки, источник спроса на жильё ---
и что происходит с выдачами, когда такой источник закрывается.

Прямой ответ на этот вопрос упирается в проблему идентификации. Доля
льготных программ в выдачах региона коррелирует и со ставкой, и с
ценами, и с доходами, поэтому регрессия выдач на льготную долю не
отделяет причинного эффекта от совместного тренда. Мы используем
событие, которое даёт внешний источник вариации: прекращение массовой
программы «Льготная ипотека 8\%» 1 июля 2024 года. Логика дизайна
состоит в том, что закрытие программы должно сильнее сказаться там, где
она занимала большую долю выдач, --- и эта доля измерена по данным,
относящимся исключительно к периоду до события.

Работа вносит три вклада. Во-первых, она показывает, что
дифференциальный эффект закрытия субсидируемой программы в России
идентифицируется, но только при достаточно широкой контрольной группе:
на панели из девяти субъектов Уральского макрорегиона дизайн не проходит
собственные проверки на состоятельность, тогда как на 84 субъектах
предпосылка параллельных трендов выполняется. Во-вторых, эффект
обнаруживается в объёме выдач, но не в их числе, что указывает на
изменение размера кредита, а не на изменение числа сделок. В-третьих,
работа документирует, что при малом числе кластеров обычные стандартные
ошибки дают завышенную значимость: первоначальная оценка с t = 5,04 при
кластеризации по региону падает до t = 1,83.

\section{Обзор
литературы}\label{ux43eux431ux437ux43eux440-ux43bux438ux442ux435ux440ux430ux442ux443ux440ux44b}

Литература о связи ДКП и рынка жилья распадается на три линии.

\textbf{Каналы трансмиссии.} Структурная рамка задана работами
Iacoviello (\citeproc{ref-iacoviello2005}{2005 г.}), который показал,
что кредитные ограничения, обеспеченные стоимостью жилья, усиливают
эффект шоков спроса, и Mishkin (\citeproc{ref-mishkin2007}{2007 г.}),
систематизировавшего шесть каналов передачи --- от стоимости капитала и
ожиданий до балансового и риск-аппетита. Taylor
(\citeproc{ref-taylor2007}{2007 г.}) на контрфактическом анализе
показал, что отклонение ставки от правила в 2002--2005 годах само стало
источником цикла на рынке жилья. Для нас важна общая мысль этой линии:
скорость и сила передачи зависят от институциональной структуры
ипотечного рынка, а не только от уровня ставки.

\textbf{Российские оценки.} Sinyakov и Shelovanova
(\citeproc{ref-sinyakov2023}{2023 г.}, \citeproc{ref-sinyakov2025}{2025
г.}) на микроданных обследования домохозяйств показали слабую
эластичность спроса по ставке: рост ставки на 1 п.п. снижает вероятность
обращения за кредитом лишь на 1,5--2,3 \%, а заметный эффект появляется
при изменениях порядка 10 п.п. Demidova и Shchankina
(\citeproc{ref-demidova2025}{2025 г.}) на панели 85 регионов обнаружили,
что связь ключевой ставки с ипотечными ставками распалась после 2020
года и что субсидируемые программы ослабляют передачу. Zvereva
(\citeproc{ref-zvereva2025}{2025 г.}) документирует региональную
асимметрию передачи и пространственные эффекты. Smirnova
(\citeproc{ref-smirnova2025}{2025 г.}) обращает внимание на двойственную
функцию жилья --- потребительского блага и инвестиционного актива, ---
из-за которой низкие ставки стимулируют обе функции одновременно.

\textbf{Макропруденциальные инструменты.} Kuttner и Shim
(\citeproc{ref-kuttner2013}{2013 г.}) на панели 57 экономик показали,
что ограничения долговой нагрузки работают против роста ипотечного
кредитования, а налоги на жильё --- против цен. Лаптева
(\citeproc{ref-lapteva2025}{2025 г.}) на данных 591 банка подтверждает
результативность макропруденциальных мер в России с лагом около двух
кварталов и асимметрией ужесточения над смягчением. Обзор политики Банка
России (\citeproc{ref-cbr2025mpp}{2025 г.}) фиксирует, что
макропруденциальные лимиты радикально изменили структуру выдач: доля
кредитов с показателем долговой нагрузки выше 80 \% в долевом
строительстве снизилась с 46 \% до 3 \%.

Настоящая работа продолжает вторую линию, но меняет источник вариации:
вместо временных рядов ставок используется дискретное институциональное
событие и пространственная вариация экспозиции к нему. В этом смысле она
ближе к логике событийных исследований, применяемых в
макропруденциальной литературе (\citeproc{ref-kuttner2013}{Kuttner и
Shim 2013 г.}).

\section{Институциональный
контекст}\label{ux438ux43dux441ux442ux438ux442ux443ux446ux438ux43eux43dux430ux43bux44cux43dux44bux439-ux43aux43eux43dux442ux435ux43aux441ux442}

Массовая программа «Льготная ипотека 8\%» действовала с 2020 года и была
прекращена 1 июля 2024 года. Помимо неё на рынке оставались адресные
программы: «Семейная ипотека», ИТ-ипотека и «Дальневосточная и
арктическая ипотека». Ключевое для дизайна обстоятельство состоит в том,
что массовая программа была наиболее широкой по кругу заёмщиков и не
имела ограничений, привязанных к составу семьи или профессии, --- именно
поэтому её закрытие представляет собой шок, сила которого различается
между регионами пропорционально прежней роли программы.

С 1 октября 2026 года вступает в силу реформа «Семейной ипотеки»: ставка
дифференцируется по числу детей и возрасту младшего ребёнка, лимиты
меняются, срок субсидирования ограничивается. На дату расчёта данные по
выдачам заканчиваются сентябрём 2026 года, поэтому эффект этой реформы в
работе не оценивается; экспозиция к ней рассчитана и приводится как
описательная величина для последующих оценок.

\section{Данные}\label{ux434ux430ux43dux43dux44bux435}

\textbf{Выдачи жилищных кредитов.} Источник --- статистика Банка России
по субъектам Российской Федерации (формы 02\_06 и 02\_07), помесячно с
января 2019 года по август 2026 года, в двух метриках: число выданных
кредитов и объём в рублях (\citeproc{ref-cbrdata}{Банк России 2026 г.}).
Панель охватывает 84 единицы территориального деления: все субъекты
федерации, а также Москву, Санкт-Петербург и Севастополь, которые в
исходном файле находятся на отдельном уровне агрегации.

\textbf{Программные выдачи.} Источник --- открытые отчёты ДОМ.РФ по
программам с государственной поддержкой в региональном разрезе
(\citeproc{ref-domrfdata}{ДОМ.РФ 2026 г.}). Ряды по массовой программе
ведутся как накопленный итог, поэтому помесячный поток восстановлен
разностью последних снимков внутри месяца.

\textbf{Мера экспозиции.} Экспозиция региона определяется как доля выдач
по массовой программе в общих жилищных выдачах за двенадцать месяцев до
события, то есть с июля 2023 года по июнь 2024 года. Все компоненты этой
меры относятся к периоду до события, что исключает механическую связь с
последующей динамикой.

\textbf{Контрольные и описательные ряды.} Цены на первичном и вторичном
рынке жилья по субъектам (квартальные средние, Росстат), демографическая
структура домохозяйств по переписи 2020 года, численность детей в
возрасте до шести лет. Эти ряды используются как предопределённые
контроли в спецификациях для уральской панели и как описательные
характеристики выборки.

\textbf{Дефекты источников, учтённые при сборке.} Три особенности
исходных данных требуют явной фиксации. Первая: строка «в том числе
Ненецкий автономный округ» является вложенной внутри Архангельской
области, и её включение даёт двойной счёт. Вторая: регион «Тюменская
область (без Югры и ЯНАО)» присутствует в файле дважды с разными
федеральными округами --- значения 1 461 и 3 927; корректной является
уральская строка, поскольку только она сходится с целым (3 927 + 3 848 +
1 421 = 9 196). Третья: ДОМ.РФ публикует Тюменскую область целиком,
включая автономные округа, поэтому сопоставимый ряд для собственно
Тюменской области получен вычитанием. Это наследует известный дефект:
выдачи автономий в отчётах ДОМ.РФ занижены.

Экспозиция измерена для 84 территориальных единиц; медиана 0.161,
минимум 0.000, максимум 0.490.

Описательные статистики меры экспозиции приведены в таблице
\ref{tab:summary}, десять наиболее экспонированных регионов --- в
таблице \ref{tab:exposure}, распределение --- на рисунке
\ref{fig:exposure}.

\begin{table}[htbp]
\centering
\caption{Экспозиция регионов к закрытию массовой программы}
\label{tab:summary}
\begin{tabular}{lr}
\toprule
Показатель & Значение \\
\midrule
Число единиц & 84 \\
Среднее & 0.166 \\
Медиана & 0.161 \\
Стандартное отклонение & 0.096 \\
Минимум & 0 \\
Первый дециль & 0.041 \\
Девятый дециль & 0.275 \\
Максимум & 0.49 \\
\bottomrule
\end{tabular}
\par\smallskip\begin{minipage}{0.92\textwidth}\footnotesize
Экспозиция — доля выдач по программе «Льготная 8%» в общих жилищных выдачах за 12 месяцев до события (июль 2023 — июнь 2024). 
\end{minipage}
\end{table}

\begin{table}[htbp]
\centering
\caption{Десять регионов с наибольшей экспозицией}
\label{tab:exposure}
\begin{tabular}{lr}
\toprule
Регион & Экспозиция \\
\midrule
Ленинградская область & 0.4902 \\
Краснодарский край & 0.4227 \\
Город Санкт-Петербург & 0.375 \\
Город Москва & 0.3584 \\
Калининградская область & 0.3488 \\
Московская область & 0.333 \\
Воронежская область & 0.2975 \\
Новосибирская область & 0.2883 \\
Свердловская область & 0.2775 \\
Республика Хакасия & 0.2703 \\
Тульская область & 0.27 \\
Рязанская область & 0.2643 \\
Республика Адыгея & 0.2632 \\
Ростовская область & 0.2564 \\
Удмуртская Республика & 0.248 \\
Пензенская область & 0.2476 \\
Республика Тыва & 0.2362 \\
Республика Татарстан & 0.2265 \\
Ярославская область & 0.2208 \\
Вологодская область & 0.2054 \\
Ставропольский край & 0.2037 \\
Красноярский край & 0.2034 \\
Республика Башкортостан & 0.2013 \\
Ульяновская область & 0.1998 \\
Иркутская область & 0.1966 \\
Пермский край & 0.1966 \\
Курская область & 0.1916 \\
Тверская область & 0.19 \\
Калужская область & 0.1859 \\
Алтайский край & 0.1853 \\
Самарская область & 0.182 \\
Нижегородская область & 0.1793 \\
Астраханская область & 0.1787 \\
Волгоградская область & 0.1759 \\
Орловская область & 0.1759 \\
Псковская область & 0.1741 \\
Новгородская область & 0.1729 \\
Брянская область & 0.1714 \\
Республика Карелия & 0.1704 \\
Карачаево-Черкесская Республика & 0.1633 \\
Челябинская область & 0.1629 \\
Чувашская Республика & 0.1621 \\
Республика Марий Эл & 0.1605 \\
Смоленская область & 0.1582 \\
Липецкая область & 0.1581 \\
Костромская область & 0.1577 \\
Томская область & 0.1553 \\
Республика Крым & 0.1522 \\
Владимирская область & 0.1489 \\
Кировская область & 0.1489 \\
Оренбургская область & 0.1459 \\
Белгородская область & 0.1457 \\
Республика Мордовия & 0.1422 \\
Приморский край & 0.137 \\
Курганская область & 0.1348 \\
Ивановская область & 0.1346 \\
Республика Северная Осетия & 0.1341 \\
Кемеровская область & 0.1294 \\
Тамбовская область & 0.1293 \\
Омская область & 0.1267 \\
Архангельская область & 0.1263 \\
Сахалинская область & 0.112 \\
Республика Коми & 0.1083 \\
Город Севастополь & 0.1079 \\
Хабаровский край & 0.1046 \\
Саратовская область & 0.1035 \\
Амурская область & 0.0999 \\
Камчатский край & 0.0908 \\
Республика Бурятия & 0.0825 \\
Республика Саха & 0.0807 \\
Республика Калмыкия & 0.0784 \\
Забайкальский край & 0.0736 \\
Магаданская область & 0.0699 \\
Кабардино-Балкарская Республика & 0.0608 \\
ХМАО-Югра & 0.0595 \\
ЯНАО & 0.0332 \\
Республика Алтай & 0.0215 \\
Мурманская область & 0.0192 \\
Республика Дагестан & 0.014 \\
Чеченская Республика & 0.0111 \\
Еврейская АО & 0 \\
Республика Ингушетия & 0 \\
Тюменская область & 0 \\
Чукотский АО & 0 \\
\bottomrule
\end{tabular}
\par\smallskip\begin{minipage}{0.92\textwidth}\footnotesize
Полный список из 84 единиц — в файле \texttt{tables/tab\_exposure.csv}. 
\end{minipage}
\end{table}

\begin{figure}[htbp]
\centering
\includegraphics[width=0.78\textwidth]{figures/fig_exposure.pdf}
\caption{Распределение экспозиции регионов к закрытию массовой программы}
\label{fig:exposure}
\end{figure}

\section{Методология}\label{ux43cux435ux442ux43eux434ux43eux43bux43eux433ux438ux44f}

\subsection{Оцениваемая
величина}\label{ux43eux446ux435ux43dux438ux432ux430ux435ux43cux430ux44f-ux432ux435ux43bux438ux447ux438ux43dux430}

Обозначим через \(y_{rt}\) логарифм показателя выдач в регионе \(r\) в
месяц \(t\), через \(E_r\) --- экспозицию региона, измеренную до
события, и через \(D_t\) --- индикатор периода после 1 июля 2024 года.
Оценивается уравнение

\[ y_{rt} = \alpha_r + \lambda_t + \beta\,(D_t \times E_r) + \gamma_r t + \varepsilon_{rt}, \]

где \(\alpha_r\) и \(\lambda_t\) --- фиксированные эффекты региона и
месяца, а \(\gamma_r t\) --- региональный линейный тренд, включаемый в
расширенных спецификациях. Коэффициент \(\beta\) имеет смысл
дифференциального эффекта на единицу экспозиции. Поскольку \(E_r\)
постоянна внутри региона, её уровень поглощается фиксированным эффектом
региона, а общий макрошок --- фиксированным эффектом месяца;
идентификация опирается исключительно на дифференциальную реакцию.

\subsection{Событийное
исследование}\label{ux441ux43eux431ux44bux442ux438ux439ux43dux43eux435-ux438ux441ux441ux43bux435ux434ux43eux432ux430ux43dux438ux435}

Для проверки предпосылки параллельных трендов оценивается динамическая
спецификация, в которой экспозиция взаимодействуется с дамми кварталов
относительно события, а базовым выступает последний предсобытийный
квартал. Незначимость предсобытийных коэффициентов является условием
содержательной интерпретации \(\beta\).

\subsection{Вывод при малом числе
кластеров}\label{ux432ux44bux432ux43eux434-ux43fux440ux438-ux43cux430ux43bux43eux43c-ux447ux438ux441ux43bux435-ux43aux43bux430ux441ux442ux435ux440ux43eux432}

Стандартные ошибки кластеризуются по региону. Поскольку число кластеров
ограничено, дополнительно применяется рандомизационный тест: экспозиция
случайно перераспределяется между регионами, спецификация
переоценивается, и распределение полученных коэффициентов сравнивается с
наблюдённым (\citeproc{ref-young2019}{Young 2019 г.}). Такая процедура
не опирается на асимптотику по числу кластеров и устойчива к малым
выборкам (\citeproc{ref-cameron2008}{Cameron и др. 2008 г.}).

\subsection{Плацебо-тест}\label{ux43fux43bux430ux446ux435ux431ux43e-ux442ux435ux441ux442}

В качестве проверки на ложную идентификацию оценивается та же
спецификация с фиктивным событием 1 июля 2023 года на предсобытийном
окне. Содержательное требование состоит в том, чтобы эффект на фиктивном
событии был незначим и не воспроизводил знак и величину основного
эффекта.

\section{Результаты}\label{ux440ux435ux437ux443ux43bux44cux442ux430ux442ux44b}

\subsection{Основные оценки по всем
субъектам}\label{ux43eux441ux43dux43eux432ux43dux44bux435-ux43eux446ux435ux43dux43aux438-ux43fux43e-ux432ux441ux435ux43c-ux441ux443ux431ux44aux435ux43aux442ux430ux43c}

\begin{table}[htbp]
\centering
\caption{Закрытие массовой программы: число ипотечных кредитов}
\label{tab:count}
\begin{tabular}{lrrrr}
\toprule
Спецификация & Коэффициент & Ст. ошибка & N & Кластеров \\
\midrule
базовая TWFE & -0.112 & (0.175) & 2016 & 84 \\
+ региональные тренды & -0.067 & (0.279) & 2016 & 84 \\
+ предразмер × пост & +0.059 & (0.223) & 2016 & 84 \\
+ тренды + предразмер & -0.213 & (0.311) & 2016 & 84 \\
+ тренды + предразмер, взвеш. & -0.453* & (0.243) & 2016 & 84 \\
\bottomrule
\end{tabular}
\par\smallskip\begin{minipage}{0.92\textwidth}\footnotesize
Зависимая переменная — логарифм числа выданных кредитов. Окно: июль 2023 — июнь 2025. Во всех спецификациях — фиксированные эффекты региона и месяца; стандартные ошибки кластеризованы по региону. $^{*}p<0{,}10$, $^{**}p<0{,}05$, $^{***}p<0{,}01$. 
\end{minipage}
\end{table}

\begin{table}[htbp]
\centering
\caption{Закрытие массовой программы: объём ипотечных выдач}
\label{tab:volume}
\begin{tabular}{lrrrr}
\toprule
Спецификация & Коэффициент & Ст. ошибка & N & Кластеров \\
\midrule
базовая TWFE & -0.118 & (0.169) & 2016 & 84 \\
+ региональные тренды & -0.557*** & (0.145) & 2016 & 84 \\
+ предразмер × пост & -0.004 & (0.221) & 2016 & 84 \\
+ тренды + предразмер & -0.582*** & (0.213) & 2016 & 84 \\
+ тренды + предразмер, взвеш. & -0.442** & (0.201) & 2016 & 84 \\
\bottomrule
\end{tabular}
\par\smallskip\begin{minipage}{0.92\textwidth}\footnotesize
Зависимая переменная — логарифм объёма выдач в млн рублей. Обозначения те же, что в таблице \ref{tab:count}. 
\end{minipage}
\end{table}

По числу кредитов эффект не обнаруживается ни в одной спецификации:
коэффициенты колеблются от \(-0{,}112\) до \(+0{,}059\) при кластерных
стандартных ошибках, исключающих значимость (таблица \ref{tab:count}).
По объёму выдач картина иная: с региональными трендами коэффициент равен
\(-0{,}557\) при \(t = -3{,}84\), с добавлением предразмерного контроля
\(-0{,}582\) при \(t = -2{,}73\), при взвешивании по размеру региона
\(-0{,}442\) при \(t = -2{,}20\) (таблица \ref{tab:volume}).
Рандомизационный тест для спецификации по объёму с трендами даёт
\(p = 0{,}003\), тогда как для числа кредитов он равен \(0{,}466\) в
базовой спецификации и \(0{,}335\) в спецификации с трендами.

Содержательно расхождение между числом и объёмом означает, что закрытие
массовой программы срезало прежде всего размер кредита: уходил крупный
субсидируемый кредит, который замещался меньшим по величине рыночным.
Это согласуется с конструкцией программы, действовавшей с лимитом 6 млн
рублей и занимавшей наибольшую долю именно в регионах с дорогим жильём.

\subsection{Динамика
эффекта}\label{ux434ux438ux43dux430ux43cux438ux43aux430-ux44dux444ux444ux435ux43aux442ux430}

\begin{table}[htbp]
\centering
\caption{Событийное исследование: динамика эффекта}
\label{tab:event}
\begin{tabular}{lrr}
\toprule
Квартал к событию & Коэффициент & Ст. ошибка \\
\midrule
-4 & +0.419 & (0.336) \\
-3 & +0.291 & (0.181) \\
-2 & -0.260 & (0.222) \\
0 & -0.518** & (0.205) \\
1 & -0.715*** & (0.228) \\
2 & -1.026** & (0.404) \\
3 & -1.593*** & (0.409) \\
\bottomrule
\end{tabular}
\par\smallskip\begin{minipage}{0.92\textwidth}\footnotesize
Зависимая переменная — логарифм числа ипотечных кредитов. Каждый коэффициент — взаимодействие экспозиции с дамми квартала относительно события 01.07.2024; базовый квартал $-1$ опущен. Включены фиксированные эффекты региона и месяца и региональные линейные тренды. 
\end{minipage}
\end{table}

\begin{figure}[htbp]
\centering
\includegraphics[width=0.82\textwidth]{figures/fig_event_study.pdf}
\caption{Динамика эффекта закрытия массовой программы (событийное исследование)}
\label{fig:event}
\end{figure}

Предсобытийные коэффициенты незначимы, что поддерживает предпосылку
параллельных трендов: оценки составляют \(+0{,}419\) (\(t = 1{,}25\)),
\(+0{,}291\) (\(t = 1{,}61\)) и \(-0{,}260\) (\(t = -1{,}17\)) для
кварталов \(-4\), \(-3\) и \(-2\) соответственно. Постсобытийные
коэффициенты значимы и монотонно растут по модулю: \(-0{,}518\)
(\(t = -2{,}52\)) в квартале события, \(-0{,}715\) (\(t = -3{,}14\)),
\(-1{,}026\) (\(t = -2{,}54\)) и \(-1{,}593\) (\(t = -3{,}90\)) на
третьем квартале после события (таблица \ref{tab:event}, рисунок
\ref{fig:event}). В пересчёте на регионы эффект на третьем квартале
составляет \(-22{,}6\) \% при медианной экспозиции \(0{,}161\),
\(-35{,}5\) \% на девятом дециле и \(-54{,}2\) \% в Ленинградской
области. Средний эффект по четырём постсобытийным кварталам при
медианной экспозиции близок к \(-14\) \%.

\subsection{Проверка
плацебо}\label{ux43fux440ux43eux432ux435ux440ux43aux430-ux43fux43bux430ux446ux435ux431ux43e}

Фиктивное событие 1 июля 2023 года даёт коэффициент \(+0{,}381\) при
\(t = 1{,}81\): знак противоположен основному эффекту, величина
незначима. Иными словами, предсобытийная динамика не воспроизводит ни
направления, ни масштаба оценки, полученной на настоящем событии.

\subsection{Контраст: уральская
панель}\label{ux43aux43eux43dux442ux440ux430ux441ux442-ux443ux440ux430ux43bux44cux441ux43aux430ux44f-ux43fux430ux43dux435ux43bux44c}

\begin{table}[htbp]
\centering
\caption{Уральская панель (9 регионов): эффект не идентифицируется}
\label{tab:ural}
\begin{tabular}{lrrrrr}
\toprule
Зависимая & Спецификация & Коэффициент & Ст. ошибка & N & Кластеров \\
\midrule
Общие выдачи & базовая & +0.064 & (0.071) & 216 & 9 \\
Общие выдачи & + контроли × пост & -0.046 & (0.134) & 216 & 9 \\
Общие выдачи & + тренды + контроли & -0.078 & (0.257) & 216 & 9 \\
Рыночные & базовая & +0.250 & (0.236) & 216 & 9 \\
Рыночные & + контроли × пост & +0.228 & (0.182) & 216 & 9 \\
Рыночные & + тренды + контроли & +1.249*** & (0.279) & 216 & 9 \\
Льготные без массовой & базовая & +2.059* & (1.127) & 214 & 9 \\
Льготные без массовой & + контроли × пост & +0.965** & (0.479) & 214 & 9 \\
Льготные без массовой & + тренды + контроли & +3.698*** & (0.232) & 214 & 9 \\
\bottomrule
\end{tabular}
\par\smallskip\begin{minipage}{0.92\textwidth}\footnotesize
Выборка ограничена девятью субъектами Уральского макрорегиона. Контроли — население, доля детей 0--6 лет, структура семей по ВПН-2020, цена первичного рынка 2023 года. 
\end{minipage}
\end{table}

Тот же дизайн на панели из девяти субъектов Уральского макрорегиона не
идентифицируется (таблица \ref{tab:ural}). По общим выдачам коэффициент
меняет знак от спецификации к спецификации: \(+0{,}064\) в базовой,
\(-0{,}046\) с контролями, \(-0{,}078\) с трендами и контролями.
Рандомизационный тест даёт \(p = 0{,}655\), плацебо-событие значимо
(\(+0{,}550\) при \(t = 6{,}13\)), предсобытийные коэффициенты отличны
от нуля. Расширение контрольной группы до всех субъектов устраняет все
три проблемы одновременно, что указывает на источник слабости: не набор
контролей, а вырожденность распределения экспозиции при девяти
наблюдениях, два из которых имеют нулевую измеренную экспозицию из-за
дефекта отчётности ДОМ.РФ по автономным округам.

\section{Робастность}\label{ux440ux43eux431ux430ux441ux442ux43dux43eux441ux442ux44c}

\textbf{Окно измерения экспозиции.} Замена двенадцатимесячного окна на
шести- и восемнадцатимесячное не меняет выводов: коэффициенты по числу
кредитов равны \(-0{,}124\), \(-0{,}112\) и \(-0{,}112\) соответственно.

\textbf{Взвешивание.} Взвешивание наблюдений по среднему размеру региона
увеличивает абсолютную величину коэффициента по объёму до \(-0{,}442\),
значимость сохраняется на пятипроцентном уровне.

\textbf{Состав контрольной группы.} Исключение Тюменской области, для
которой экспозиция получена вычитанием и потому наследует дефект
отчётности, входит в перечень ближайших проверок; это единственный
регион с такой конструкцией ряда.

\section{Обсуждение}\label{ux43eux431ux441ux443ux436ux434ux435ux43dux438ux435}

Полученные оценки укладываются в рамку ослабленной трансмиссии,
описанную в российской литературе (\citeproc{ref-sinyakov2023}{Sinyakov
и Shelovanova 2023 г.}; \citeproc{ref-demidova2025}{Demidova и
Shchankina 2025 г.}). Если программы с субсидируемой ставкой формируют
автономный спрос, то их закрытие должно вызывать сокращение выдач,
пропорциональное прежней роли программы, --- именно это и показывают
оценки. Тот факт, что эффект обнаруживается в объёме, но не в числе
кредитов, добавляет к этой картине важную деталь: субсидия поддерживала
не количество сделок, а их масштаб. Для политики это означает, что
отмена программы сказывается прежде всего на структуре кредитования ---
средний размер кредита снижается, --- а не на числе домохозяйств,
получивших кредит.

Наблюдаемая динамика --- нарастание эффекта от квартала события к
третьему кварталу --- согласуется с постепенным исчерпанием ранее
одобренных заявок и с инерцией адаптации заёмщиков к рыночным ставкам.
Мгновенного обрыва выдач не происходит, что важно учитывать при
интерпретации месячных данных вокруг дат изменения программ.

\section{Ограничения}\label{ux43eux433ux440ux430ux43dux438ux447ux435ux43dux438ux44f}

Первое ограничение связано с конструкцией ряда экспозиции для Тюменской
области без автономных округов: значение получено вычитанием, поскольку
источник публикует область целиком. Второе: объёмный канал не разложен
на число сделок и средний размер кредита, из-за чего интерпретация через
размер кредита остаётся содержательной гипотезой, а не измеренным
фактом. Третье: эффект идентифицируется дифференциально, то есть общий
макрошок, действующий на все регионы одинаково, поглощается
фиксированными эффектами месяца и в оценку не входит; работа отвечает на
вопрос о различиях между регионами, а не о совокупном эффекте программы.
Четвёртое: реформа «Семейной ипотеки» от 1 октября 2026 года на
имеющихся данных не наблюдаема. Наконец, оценки не отделяют эффект
закрытия программы от сопутствующих изменений макропруденциальной
политики того же периода.

\section{Заключение}\label{ux437ux430ux43aux43bux44eux447ux435ux43dux438ux435}

Закрытие массовой программы льготной ипотеки 1 июля 2024 года даёт
редкую возможность оценить, насколько программы с субсидируемой ставкой
формировали автономный спрос на жильё. На панели из 84 субъектов
Российской Федерации дифференциальный эффект на объём выдач отрицателен,
статистически значим и нарастает во времени; предпосылка параллельных
трендов выполняется, а плацебо-событие эффект не воспроизводит. Эффект
выражен в объёме, но не в числе кредитов. Сопоставление с ограниченной
уральской панелью показывает, что решающим условием корректной
идентификации в такой задаче выступает широта контрольной группы, а не
перечень контролей: при девяти регионах дизайн не проходит собственные
проверки, при восьмидесяти четырёх --- проходит.

\section{Приложение А. Механика формирования статистики инвестиций в
основной
капитал}\label{ux43fux440ux438ux43bux43eux436ux435ux43dux438ux435-ux430.-ux43cux435ux445ux430ux43dux438ux43aux430-ux444ux43eux440ux43cux438ux440ux43eux432ux430ux43dux438ux44f-ux441ux442ux430ux442ux438ux441ux442ux438ux43aux438-ux438ux43dux432ux435ux441ux442ux438ux446ux438ux439-ux432-ux43eux441ux43dux43eux432ux43dux43eux439-ux43aux430ux43fux438ux442ux430ux43b}

Оценки спроса на жильё, полученные по выдачам кредитов, в дальнейшем
должны быть связаны с показателями строительной активности. Для этого
необходимо понимать, как устроена статистика инвестиций в основной
капитал (ИОК), поскольку она устроена иначе, чем статистика выдач.

ИОК формируется не в момент ввода объекта в эксплуатацию, а по мере
признания затрат в бухгалтерском учёте: строительно-монтажные работы
учитываются в размере фактически выполненного объёма по справке о
стоимости работ независимо от момента оплаты, оборудование --- после
поступления на место назначения и оприходования заказчиком. Прямой круг
наблюдения образуют шесть форм: квартальная форма П-2 для юридических
лиц, не относящихся к субъектам малого предпринимательства (срок ---
20-е число после отчётного периода, за январь--декабрь --- 8 февраля);
годовая форма П-2 (инвест) (1 апреля); выборочные формы ПМ и МП (микро)
для малых и микропредприятий; форма 1-ВЭС для предприятий с иностранным
участием; форма ДАФЛ для лизинговых организаций. Не включаются в ИОК
ремонт (как текущий, так и капитальный), основные средства дешевле 40
тыс. рублей за единицу, бывшие в употреблении объекты, квартиры и
непроизведённые активы --- земля, аренда, лицензии
(\citeproc{ref-rosstat2026iok}{Росстат 2016 г.}).

Внутри года полный круг получается умножением данных крупных и средних
организаций на коэффициент досчёта, рассчитанный по окончательным итогам
предыдущего года; за год формируется пять последовательных оценок,
последние из которых опираются на балансовый метод --- восстановление
объёма инвестиций по ресурсам строительных материалов и машин с
корректировками на экспорт, ремонт и конечное потребление домашних
хозяйств (\citeproc{ref-rosstat2014reg}{Росстат 2014 г.}). Практическое
следствие для нашего исследования состоит в том, что ИОК публикуется
квартально, содержит досчитанную ненаблюдаемую часть и пересматривается,
поэтому сопоставление выдач с инвестициями требует согласования не
только по периоду, но и по моменту признания.

\section{Приложение Б.
Воспроизводимость}\label{ux43fux440ux438ux43bux43eux436ux435ux43dux438ux435-ux431.-ux432ux43eux441ux43fux440ux43eux438ux437ux432ux43eux434ux438ux43cux43eux441ux442ux44c}

Все расчёты воспроизводятся скриптами \texttt{did\_all\_regions.py}
(оценки по всем субъектам), \texttt{ufo\_did\_controls.py} (уральская
панель) и \texttt{export\_for\_paper.py} (таблицы и рисунки). Исходные
данные --- публичные ряды Банка России и ДОМ.РФ, загруженные в базу
исследовательского трека. Сборка документа: \texttt{make\ pdf} либо
\texttt{Rscript\ -e\ \textquotesingle{}rmarkdown::render("paper.Rmd")\textquotesingle{}}.

\appendix

\protect\phantomsection\label{refs}
\begin{CSLReferences}{1}{1}
\bibitem[\citeproctext]{ref-cameron2008}
Cameron, A. Colin, Jonah B. Gelbach, и Douglas L. Miller. 2008 г.
{«Bootstrap-Based Improvements for Inference with Clustered Errors»}.
\emph{Review of Economics and Statistics} 90 (3): 414--27.
\url{https://doi.org/10.1162/rest.90.3.414}.

\bibitem[\citeproctext]{ref-demidova2025}
Demidova, Olga A., и Anastasia A. Shchankina. 2025 г. {«Impact of Key
Interest Rate Changes on Mortgage Rates in Russian Regions»}.
\emph{Экономика и математические методы} 61 (2): 75--89.
\url{https://doi.org/10.31857/S0424738825020061}.

\bibitem[\citeproctext]{ref-iacoviello2005}
Iacoviello, Matteo. 2005 г. {«House Prices, Borrowing Constraints, and
Monetary Policy in the Business Cycle»}. \emph{American Economic Review}
95 (3): 739--64. \url{https://doi.org/10.1257/0002828054201477}.

\bibitem[\citeproctext]{ref-kuttner2013}
Kuttner, Kenneth N., и Ilhyock Shim. 2013 г. \emph{Can Non-Interest Rate
Policies Stabilise Housing Markets? Evidence from a Panel of 57
Economies}. BIS Working Papers No. 433. Bank for International
Settlements.

\bibitem[\citeproctext]{ref-mishkin2007}
Mishkin, Frederic S. 2007 г. \emph{Housing and the Monetary Transmission
Mechanism}. Jackson Hole Symposium Paper. Federal Reserve Bank of Kansas
City.

\bibitem[\citeproctext]{ref-sinyakov2023}
Sinyakov, Andrey, и Tatiana Shelovanova. 2023 г. \emph{Demand for
Consumer Loans in Russia: How Strong is the Interest Rate Channel of
Monetary Policy?} Working Paper Series No. 120. Bank of Russia.

\bibitem[\citeproctext]{ref-sinyakov2025}
Sinyakov, Andrey, и Tatiana Shelovanova. 2025 г. {«Demand for Consumer
Loans in Russia: How Strong is the Interest Rate Channel of Monetary
Policy?»} \emph{Russian Journal of Economics} 11 (1): 47--75.
\url{https://doi.org/10.32609/j.ruje.11.145314}.

\bibitem[\citeproctext]{ref-smirnova2025}
Smirnova, Z. 2025 г. \emph{Real Estate Market and Monetary Policy:
Searching for Balance}. Econs.online (мнение, не отражает позицию Банка
России).

\bibitem[\citeproctext]{ref-taylor2007}
Taylor, John B. 2007 г. \emph{Housing and Monetary Policy}. NBER Working
Paper No. 13682. National Bureau of Economic Research.
\url{https://doi.org/10.3386/w13682}.

\bibitem[\citeproctext]{ref-young2019}
Young, Alwyn. 2019 г. {«Channeling Fisher: Randomization Tests and the
Statistical Insignificance of Seemingly Significant Experimental
Results»}. \emph{Quarterly Journal of Economics} 134 (2): 557--98.
\url{https://doi.org/10.1093/qje/qjy029}.

\bibitem[\citeproctext]{ref-zvereva2025}
Zvereva, V. 2025 г. {«Monetary Policy Transmission Mechanism: The Role
of Household Heterogeneity and Spatial Effects»}. \emph{Russian Journal
of Macroeconomics and Financial Economics} 17 (3).

\bibitem[\citeproctext]{ref-cbr2025mpp}
Банк России. 2025 г. \emph{Банк России принял ряд решений по
макропруденциальной политике}. Пресс-релиз.

\bibitem[\citeproctext]{ref-cbrdata}
Банк России. 2026 г. \emph{Сведения о размещённых и привлечённых
средствах: жилищные кредиты по субъектам Российской Федерации (формы
02\_06, 02\_07, 02\_13)}. Статистические данные Банка России.

\bibitem[\citeproctext]{ref-domrfdata}
ДОМ.РФ. 2026 г. \emph{Отчёты о выдаче ипотечных кредитов по программам с
государственной поддержкой (weeklyreport, региональный разрез)}.
Открытые данные АО «ДОМ.РФ».

\bibitem[\citeproctext]{ref-lapteva2025}
Лаптева, Е. В. 2025 г. {«Оценка влияния макропруденциальной политики на
сдерживание потребительского кредитования в России»}.
\emph{Экономический журнал Высшей школы экономики} 29 (4): 609--42.
\url{https://doi.org/10.17323/1813-8691-2025-29-4-609-642}.

\bibitem[\citeproctext]{ref-rosstat2014reg}
Росстат. 2014 г. \emph{Официальная статистическая методология
определения инвестиций в основной капитал на региональном уровне}.

\bibitem[\citeproctext]{ref-rosstat2026iok}
Росстат. 2016 г. \emph{Официальная статистическая методология
определения инвестиций в основной капитал на федеральном уровне}.

\end{CSLReferences}

\end{document}
